% NI64--0C
% Insult
% 14.0
% 2013-11-30

.. _m-insult:


Taktik \Ca{Vitale Bedrohung!}
Zeitkritisch (\enquote{Time is brain.}}):
Vitalfunktionen sichern, rascher Transport an eine geeignete Abteilung.}

Grundsätzlich
ist ein Schlaganfall-Patient aufgrund seiner Diagnose als **vital
bedroht** und gleichzeitig als **zeitkritisch**
(\enquote{Time is brain.}}) einzustufen. 
Eine definitive Therapie in einer Spezialabteilung
(\I{Stroke Unit}) kann andere, noch nicht abgestorbene
Teile des Hirns vor weiterem Schaden bewahren, d. h. es ist
sehr wichtig, dass diese Therapie möglichst schnell
erfolgt.
% 
% Bei der Behandlung des Patienten mit V. a. Schlaganfall
% gelten folgende Besonderheiten:


\begin{itemize}
  -   |TxMassVitMKBes|
  \begin{itemize}
    -   **Lagerung**: *leicht erhöhter Oberkörper*
    -   **Sauerstoff**: Gabe je nach SpO₂,
    unterer Grenzwert 94%,
    darüber ist von einer O₂-Gabe abzusehen, da es
    sonst zu einer Gefäßengstellung mit nachfolgender
    Mangeldurchblutung in den noch gesunden Hirnarealen
    kommen kann! \Commentary[Insult / Sauerstoffgabe]{In der originalen
    Lehrmeinungsnpassung ist der Grenzwert 95%. Im Sinne
    der Konsistenz mit der allgemeinen Lehrmeinung ist ein
    angepasster Grenzwert von 94% vertretbar
   :cite:`AsboeLvWien-2013-01-29`\ .}
    -   **Notarzt**: Grundsätzlich ist ein
    Schlaganfall-Patient aufgrund seiner Diagnose als vital
    bedroht einzustufen, er ist jedoch gleichzeitig zeitkritisch (s. o.).
    Daher soll von einer
    \EE{Nachberufung eines Notarztes \underline{abgesehen} werden,
    *wenn folgende Umstände gegeben sind*}:
    \begin{enumerate}
      -   Die Verdachtsdiagnose gilt als *sehr sicher*.
      -   Der Patient ist sonst als *stabil* (\prettyref{m:ersteinschaetzung}) eingestuft
      und
      bewusstseinsklar\footnote{Bei der Einschätzung des Bewusstseins
      darf man sich nicht durch die durch den Schlaganfall
      verursachten Symptome (Sprachstörungen, Blickstarre \ldots)
      beirren lassen.}.
      -   Es gibt sonst *keine anderen Umstände* die die Berufung
      eines Notarztes erforderlich machen.
      -   Es muss versucht werden, eine Aufnahme auf eine
      *Spezialabteilung* für Schlaganfälle zu
      organisieren\footnote{Der Erfolg ist abhängig von der
      Verfügbarkeit eines entsprechenden freien Bettes.} (Stroke Unit,
      Kontaktaufnahme mit der Leitstelle).
      -   Regionale Bestimmungen müssen beachtet werden!
    \end{enumerate}
  \end{itemize}
  -    **Blutdruck**: Blutdruckwerte bis \RRmmHg{220}{120} werden toleriert,
  wenn nicht cardiopulmonale Gründe (Myocardischämie,
  Herzinsuffizienz, \ldots) andere Grenzwerte erfordern.
  Senkung des Blutdrucks ist eine ärztliche Maßnahme.
  -   **Neurocheck** inkl. Messung des **Blutzuckers** ist
  zwingend erforderlich!
  -   **Transportentscheidung**: Wenn verfügbar: **Stroke Unit**, sonst
  Abklärung mit Leitstelle wegen Alternativen (Abt. für
  Neurologie oder Innere Medizin, mit Möglichkeit zur CT und
  Intensivüberwachung).
  Bei Anforderung einer Stroke Unit ist anzugeben:
  \begin{enumerate}
    -   Alter
    -   Beginn der Symptome
    -   Vorherige Schlaganfälle
  \end{enumerate}
\end{itemize}

\begin{Danger}
  -   Eine engmaschige Verlaufskontrolle, insbesonders des
  Bewusstseinszustandes, ist wichtig!
\end{Danger}


\cite{Aass-1.0var1,AassMk-2013var1,MA70-Protokoll-2012-10-11,AsboeLvWien-2013-01-29,Lang2012,Ma70-Ausrueckordnung-2012-01-24}
